% ==============================================================================
% CAPITOLO 5: CURVE PARAMETRICHE E INTEGRALI CURVILINEI
% ==============================================================================

\chapter{Curve Parametriche e Integrali Curvilinei}
\label{chap:curve_e_integrali_curvilinei}

Nel calcolo infinitesimale a una variabile, l'integrazione è tradizionalmente introdotta su intervalli compatti dell'asse reale $[a,b] \subset \R$. Tuttavia, la modellizzazione geometrica, fisica ed ingegneristica richiede frequentemente di misurare grandezze distribuite lungo traiettorie curvilinee arbitrarie nello spazio euclideo $\Rdue$ o $\Rtre$. Si pensi, ad esempio, al calcolo della massa o del baricentro di un filo metallico a densità variabile, alla determinazione del lavoro compiuto da un campo di forze su una particella in movimento, oppure alla quantificazione della circolazione di un fluido lungo un profilo alare.

Per estendere l'analisi a questi contesti, è necessario sviluppare una solida comprensione della \textbf{geometria differenziale locale delle curve} e costruire due distinte nozioni di integrale curvilineo: gli \textbf{integrali di prima specie} (relativi a campi scalari e invarianti rispetto al verso di percorrenza) e gli \textbf{integrali di seconda specie} (relativi a campi vettoriali e forme differenziali, sensibili all'orientazione del cammino).

% ==============================================================================
\section{Geometria Differenziale delle Curve Parametriche}
% ==============================================================================

Dal punto di vista intuitivo, una curva nello spazio è l'oggetto unidimensionale tracciato dal moto continuo di un punto materiale al variare del tempo. Tuttavia, in matematica è essenziale distinguere formalmente tra la legge oraria del moto (la \emph{curva parametrica} vera e propria) e la traccia geometrica lasciata nello spazio (il suo \emph{sostegno}).

\begin{definizione}{Curva Parametrica e Proprietà Topologiche}{def_curva_estesa}
Sia $I = [a,b] \subset \R$ un intervallo chiuso e limitato. Si definisce \textbf{curva parametrica} (o semplicemente \emph{curva}) in $\Rn$ ($n=2,3$) un'applicazione continua:
\[
\gamma: [a,b] \longrightarrow \Rn, \quad t \longmapsto \gamma(t) = \begin{pmatrix} x_1(t) \\ \vdots \\ x_n(t) \end{pmatrix}
\]
La variabile $t \in [a,b]$ è detta \textbf{parametro} (spesso interpretata fisicamente come il tempo).
\begin{itemize}
    \item L'insieme immagine $\Gamma = \gamma([a,b]) \subset \Rn$ è detto \textbf{sostegno} (o traiettoria) della curva.
    \item La curva si dice \textbf{chiusa} se il punto iniziale coincide con il punto finale: $\gamma(a) = \gamma(b)$.
    \item La curva si dice \textbf{semplice} se non ammette autointersezioni, ossia se l'applicazione $\gamma$ è iniettiva sull'intervallo semiaperto $[a,b)$. In particolare, se la curva è chiusa, l'unica coincidenza ammessa è $\gamma(a) = \gamma(b)$.
    \item La curva si dice di classe \textbf{$C^1([a,b])$} se ciascuna funzione componente $x_i: [a,b] \to \R$ è derivabile con continuità su tutto $[a,b]$ (intendendo le derivate negli estremi $a$ e $b$ come derivate unilaterali destra e sinistra).
    \item Una curva $C^1([a,b])$ si dice \textbf{regolare} se il suo vettore velocità non si annulla mai:
    \[
    \gamma'(t) = \begin{pmatrix} x_1'(t) \\ \vdots \\ x_n'(t) \end{pmatrix} \neq \bzero, \quad \forall t \in [a,b]
    \]
    \item La curva si dice \textbf{regolare a tratti} se l'intervallo $[a,b]$ può essere suddiviso in un numero finito di sottointervalli $[t_{k-1}, t_k]$ su ciascuno dei quali la restrizione di $\gamma$ è regolare.
\end{itemize}
\end{definizione}

\begin{osservazione}[Distinzione tra Curva e Sostegno]
È di fondamentale importanza non confondere la curva $\gamma$ (che è una funzione vettoriale dotata di verso e velocità di percorrenza) con il suo sostegno $\Gamma$ (che è un sottoinsieme geometrico di $\Rn$). 

Ad esempio, le due applicazioni:
\[
\gamma_1(t) = (\cos t, \sen t), \quad t \in [0, 2\pi]
\]
\[
\gamma_2(u) = (\cos(2u), \sen(2u)), \quad u \in [0, \pi]
\]
rappresentano due curve parametriche distinte (hanno domini diversi e velocità istantanea differente, poiché $\norm{\gamma_1'(t)} = 1$ mentre $\norm{\gamma_2'(u)} = 2$), ma condividono esattamente lo stesso sostegno geometrico, ossia la circonferenza unitaria $S^1 = \graffe{(x,y) \in \Rdue : x^2 + y^2 = 1}$.
\end{osservazione}

\subsection{Cinematica Locale: Velocità, Retta Tangente e Singolarità}

Se interpretiamo il parametro $t$ come il tempo e $\gamma(t)$ come il vettore posizione di un punto materiale nello spazio, la derivata prima $\gamma'(t)$ rappresenta il \textbf{vettore velocità istantanea}. La condizione di regolarità $\gamma'(t) \neq \bzero$ garantisce che il moto non subisca arresti istantanei; geometricamente, ciò assicura che in ogni punto della curva esista una e una sola \textbf{retta tangente ben definita}.

\begin{definizione}{Retta Tangente e Versore Tangente Unitario}{def_retta_tangente}
Sia $\gamma: [a,b] \to \Rn$ una curva regolare. Per ogni $t_0 \in [a,b]$, la \textbf{retta tangente} al sostegno della curva nel punto $P_0 = \gamma(t_0)$ è la retta passante per $P_0$ e parallela al vettore velocità $\gamma'(t_0)$, avente equazione parametrica:
\[
\mathbf{r}(u) = \gamma(t_0) + u \, \gamma'(t_0), \quad u \in \R
\]
Il \textbf{versore tangente unitario} $\bT(t)$ è il vettore di norma unitaria concorde con il verso del moto:
\[
\bT(t) = \frac{\gamma'(t)}{\norm{\gamma'(t)}} = \frac{1}{\sqrt{\sum_{i=1}^n [x_i'(t)]^2}} \begin{pmatrix} x_1'(t) \\ \vdots \\ x_n'(t) \end{pmatrix}
\]
\end{definizione}

\begin{center}
\begin{tikzpicture}[scale=1.15, >=Stealth]
    % Assi cartesiani 3D
    \draw[->, thick, gray!80] (0,0,0) -- (5.5,0,0) node[below left, black] {$x$};
    \draw[->, thick, gray!80] (0,0,0) -- (0,4.2,0) node[above, black] {$z$};
    \draw[->, thick, gray!80] (0,0,0) -- (0,0,5) node[below right, black] {$y$};

    % Elica cilindrica nello spazio
    \draw[ultra thick, navyblue, domain=0:3.8*pi, samples=120, variable=\t] 
        plot ({2.0*cos(\t r)}, {0.28*\t}, {2.0*sin(\t r)});

    % Punto su elica a t0 = 1.8*pi
    \coordinate (P0) at ({2.0*cos(1.8*pi r)}, {0.28*1.8*pi}, {2.0*sin(1.8*pi r)});
    \filldraw[crimsonred] (P0) circle (2.2pt) node[left=3pt, black] {$P_0 = \gamma(t_0)$};

    % Vettore velocità gamma'(t0)
    \draw[thick, crimsonred, ->, line width=1.3pt] (P0) -- ++({-2.0*sin(1.8*pi r)*0.9}, {0.28*0.9}, {2.0*cos(1.8*pi r)*0.9}) 
        node[above right=2pt, crimsonred] {$\gamma'(t_0)$ (Vettore Velocità)};

    % Versore tangente T(t0)
    \draw[thick, forestgreen, ->, line width=1.5pt] (P0) -- ++({-2.0*sin(1.8*pi r)*0.45}, {0.28*0.45}, {2.0*cos(1.8*pi r)*0.45}) 
        node[below left=2pt, forestgreen] {$\bT(t_0)$};

    % Vettore posizione r(t0) dall'origine
    \draw[thick, coolblue!80!black, ->, dashed] (0,0,0) -- (P0) node[midway, below] {$\gamma(t_0)$};

    % Retta tangente tratteggiata
    \draw[densely dotted, warmamber!90!black, thick] 
        ($(P0) - 1.2*({-2.0*sin(1.8*pi r)*0.9}, {0.28*0.9}, {2.0*cos(1.8*pi r)*0.9})$) -- 
        ($(P0) + 1.5*({-2.0*sin(1.8*pi r)*0.9}, {0.28*0.9}, {2.0*cos(1.8*pi r)*0.9})$) 
        node[right] {\small Retta tangente};

    \node[draw=navyblue, fill=subtlegray, rounded corners, align=center] at (2.8,4.0) 
        {\small Elica cilindrica: $\gamma(t) = (R\cos t, R\sen t, c t)$};
\end{tikzpicture}
\end{center}

Cosa accade se la condizione di regolarità viene a mancare? Quando $\gamma'(t_0) = \bzero$, la curva può presentare un arresto con repentino cambio di direzione, originando un \textbf{punto di cuspide} o un punto angoloso sul sostegno geometrico.

\begin{esempio}{Curva Cuspidale e Perdita di Regolarità}{es_cuspide_neil}
Si consideri la \emph{parabola semicubica} (o cuspidale di Neil) definita da:
\[
\gamma: [-1, 1] \longrightarrow \Rdue, \quad \gamma(t) = \begin{pmatrix} t^2 \\ t^3 \end{pmatrix}
\]
Calcolando il vettore derivata prima otteniamo:
\[
\gamma'(t) = \begin{pmatrix} 2t \\ 3t^2 \end{pmatrix} \implies \gamma'(0) = \begin{pmatrix} 0 \\ 0 \end{pmatrix}
\]
Nel punto $t=0$, corrispondente all'origine $(0,0)$, la velocità si annulla. Eliminando il parametro $t$ dalle equazioni cartesiane ($x = t^2 \ge 0 \implies t = \pm\sqrt{x}$, da cui $y = \pm x^{3/2}$ ovvero $y^2 = x^3$), si constata che il sostegno forma una punta acuta nell'origine con tangente orizzontale ma derivata seconda discontinua nel cambio di ramo. La curva è differenziabile dappertutto come funzione parametrica, ma non è regolare in $t=0$.
\end{esempio}

\subsection{Riparametrizzazioni e Concetto di Ascissa Curvilinea}

Dato che una medesima traiettoria può essere percorsa a velocità arbitrarie, sorge spontanea la necessità di studiare come si trasformano le grandezze geometriche al variare della parametrizzazione adottata.

\begin{definizione}{Riparametrizzazione Ammissibile}{def_riparametrizzazione}
Sia $\gamma: [a,b] \to \Rn$ una curva regolare di classe $C^1$. Si dice \textbf{cambio di parametro ammissibile} un diffeomorfismo di classe $C^1$:
\[
\phi: [c,d] \longrightarrow [a,b], \quad s \longmapsto t = \phi(s)
\]
tale che $\phi'(s) \neq 0$ per ogni $s \in [c,d]$. La curva composta:
\[
\tilde{\gamma}: [c,d] \longrightarrow \Rn, \quad \tilde{\gamma}(s) = (\gamma \circ \phi)(s) = \gamma(\phi(s))
\]
è detta \textbf{riparametrizzazione} di $\gamma$.
\begin{itemize}
    \item Se $\phi'(s) > 0$ per ogni $s \in [c,d]$, $\phi$ è strettamente crescente; in tal caso $\phi(c)=a$, $\phi(d)=b$ e la riparametrizzazione si dice \textbf{concorde} (preserva l'orientazione e il verso di percorrenza).
    \item Se $\phi'(s) < 0$ per ogni $s \in [c,d]$, $\phi$ è strettamente decrescente; in tal caso $\phi(c)=b$, $\phi(d)=a$ e la riparametrizzazione si dice \textbf{discorde} (inverte l'orientazione del percorso).
\end{itemize}
\end{definizione}

Applicando la regola di derivazione a catena (Chain Rule), il vettore velocità della curva riparametrizzata è proporzionale al vettore velocità originario:
\[
\tilde{\gamma}'(s) = \der{}{s}\left[ \gamma(\phi(s)) \right] = \gamma'(\phi(s)) \, \phi'(s)
\]
Di conseguenza, il versore tangente soddisfa:
\[
\tilde{\bT}(s) = \frac{\tilde{\gamma}'(s)}{\norm{\tilde{\gamma}'(s)}} = \frac{\gamma'(\phi(s)) \, \phi'(s)}{\norm{\gamma'(\phi(s))} \, \abs{\phi'(s)}} = \sgn(\phi'(s)) \, \bT(\phi(s))
\]
Ciò dimostra formalmente che la direzione della retta tangente è un'invariante geometrica intrinseca del sostegno, mentre il verso del versore tangente concorda o discorda a seconda del segno di $\phi'$.

\subsection{Ascissa Curvilinea come Parametro Intrinseco Naturale}

Tra tutte le possibili parametrizzazioni di una curva regolare, ne esiste una canonica e privilegiata dal punto di vista geometrico e cinematico: la parametrizzazione tramite la \textbf{lunghezza d'arco} (o \textbf{ascissa curvilinea}).

\begin{definizione}{Ascissa Curvilinea}{def_ascissa_curvilinea}
Sia $\gamma: [a,b] \to \Rn$ una curva regolare. Si definisce \textbf{ascissa curvilinea} la funzione scalare $s: [a,b] \to [0, L]$ che misura la distanza percorsa lungo la curva a partire dall'istante iniziale $a$ fino al tempo generico $t$:
\[
s(t) = \int_{a}^{t} \norm{\gamma'(u)} \du = \int_{a}^{t} \sqrt{\sum_{i=1}^n [x_i'(u)]^2} \du
\]
\end{definizione}

Per il Teorema Fondamentale del Calcolo Integrale, la funzione $s(t)$ è di classe $C^1([a,b])$ e la sua derivata temporale è:
\[
s'(t) = \der{s}{t} = \norm{\gamma'(t)} = v(t) > 0, \quad \forall t \in [a,b]
\]
Poiché la derivata $s'(t)$ è strettamente positiva su tutto l'intervallo, la funzione $s(t)$ è strettamente crescente e quindi invertibile. La sua funzione inversa $t = t(s) = s^{-1}(s)$, definita sull'intervallo $[0, L]$, è anch'essa di classe $C^1$ con derivata:
\[
\der{t}{s} = \frac{1}{s'(t(s))} = \frac{1}{\norm{\gamma'(t(s))}}
\]
Riparametrizzando la curva rispetto all'ascissa curvilinea $s$, otteniamo la curva naturale $\gamma(s) = \gamma(t(s))$. La velocità istantanea rispetto a $s$ vale:
\[
\der{\gamma}{s} = \gamma'(t(s)) \der{t}{s} = \frac{\gamma'(t(s))}{\norm{\gamma'(t(s))}} = \bT(s) \implies \norm{\der{\gamma}{s}} = \norm{\bT(s)} = 1
\]

\begin{osservazione}[Significato Fisico dell'Ascissa Curvilinea]
L'ascissa curvilinea $s$ rappresenta il \textbf{tempo proprio} registrato da un osservatore che viaggia lungo la traiettoria con velocità scalare costante e unitaria ($v = 1\text{ m/s}$). In questo sistema di riferimento naturale, la derivata rispetto a $s$ fornisce istantaneamente il versore tangente $\bT(s)$ senza necessità di normalizzazione.
\end{osservazione}

\subsection{Curvatura Geometrica e Cerchio Osculatore}

Come varia la direzione del versore tangente $\bT(s)$ man mano che ci si muove lungo la curva? La rapidità con cui la traiettoria devia da una linea retta è quantificata dalla \textbf{curvatura}.

\begin{definizione}{Curvatura e Raggio di Curvatura}{def_curvatura}
Sia $\gamma(s)$ una curva parametrizzata mediante l'ascissa curvilinea $s$, di classe $C^2$. Si definisce \textbf{curvatura scalare} $\kappa(s)$ la norma della derivata del versore tangente rispetto all'arco:
\[
\kappa(s) = \norm{\der{\bT}{s}}
\]
Se $\kappa(s) > 0$, si definiscono:
\begin{itemize}
    \item Il \textbf{versore normale principale}: $\bN(s) = \frac{1}{\kappa(s)} \der{\bT}{s}$;
    \item Il \textbf{raggio di curvatura}: $R(s) = \frac{1}{\kappa(s)}$;
    \item Il \textbf{cerchio osculatore}: la circonferenza giacente nel piano osculatore (spazio generato da $\bT$ e $\bN$), di raggio $R(s)$ e tangente alla curva nel punto $\gamma(s)$, con centro nel punto $C(s) = \gamma(s) + R(s) \bN(s)$.
\end{itemize}
\end{definizione}

Poiché $\norm{\bT(s)}^2 = \scal{\bT(s), \bT(s)} = 1$, derivando rispetto a $s$ si ottiene:
\[
\der{}{s} \scal{\bT(s), \bT(s)} = 2 \scal{\der{\bT}{s}, \bT(s)} = 0 \implies \der{\bT}{s} \perp \bT(s)
\]
Il vettore $\der{\bT}{s}$ è dunque sempre \textbf{ortogonale} al versore tangente, confermando che $\bN(s)$ è ortogonale a $\bT(s)$.

Per una generica parametrizzazione $t$ non naturale, la formula della curvatura in $\Rtre$ si esprime mediante il prodotto vettoriale:
\[
\kappa(t) = \frac{\norm{\gamma'(t) \times \gamma''(t)}}{\norm{\gamma'(t)}^3}
\]
Mentre per una curva piana $\gamma(t) = (x(t), y(t))$ in $\Rdue$, la formula assume la celebre espressione cartesiana:
\[
\kappa(t) = \frac{\abs{x'(t) y''(t) - y'(t) x''(t)}}{\left( [x'(t)]^2 + [y'(t)]^2 \right)^{3/2}}
\]

% ==============================================================================
\section{Rettificazione e Lunghezza delle Curve}
% ==============================================================================

Il problema classico della \emph{rettificazione di una curva} consiste nell'associare una lunghezza finita a un arco curvilineo continuo mediante l'approssimazione con spezzate poligonali inscritte.

\begin{definizione}{Poligonale Inscritta e Curva Rettificabile}{def_rettificabile}
Sia $\gamma: [a,b] \to \Rn$ una curva continua. Sia $\mathcal{P} = \graffe{a = t_0 < t_1 < \dots < t_N = b}$ una generica partizione dell'intervallo $[a,b]$. La lunghezza della poligonale inscritta associata alla partizione $\mathcal{P}$ è definita come la somma delle lunghezze euclidee delle corde:
\[
L(\mathcal{P}, \gamma) = \sum_{k=1}^N \norm{\gamma(t_k) - \gamma(t_{k-1})}
\]
La curva $\gamma$ si dice \textbf{rettificabile} se l'insieme delle lunghezze poligonali è superiormente limitato. In tal caso, la \textbf{lunghezza} di $\gamma$ è l'estremo superiore:
\[
L(\gamma) = \sup_{\mathcal{P}} L(\mathcal{P}, \gamma)
\]
\end{definizione}

\begin{center}
\begin{tikzpicture}[scale=1.2, >=Stealth]
    % Curva continua
    \draw[ultra thick, navyblue] plot [smooth, tension=0.8] coordinates {(0,0.5) (1.5,2.2) (3.5,2.0) (5.5,3.2) (7,1.8)};
    \node[navyblue, above] at (3.5,2.2) {Curva $\gamma(t)$};

    % Nodi della partizione
    \coordinate (P0) at (0,0.5);
    \coordinate (P1) at (1.2,1.9);
    \coordinate (P2) at (2.4,2.3);
    \coordinate (P3) at (3.8,1.95);
    \coordinate (P4) at (5.2,3.1);
    \coordinate (P5) at (7,1.8);

    % Poligonale inscritta
    \draw[thick, crimsonred, dashed] (P0) -- (P1) -- (P2) -- (P3) -- (P4) -- (P5);

    \filldraw[black] (P0) circle (2pt) node[below left] {$\gamma(t_0) = \gamma(a)$};
    \filldraw[black] (P1) circle (1.8pt) node[above left] {$\gamma(t_1)$};
    \filldraw[black] (P2) circle (1.8pt) node[above] {$\gamma(t_2)$};
    \filldraw[black] (P3) circle (1.8pt) node[below] {$\gamma(t_3)$};
    \filldraw[black] (P4) circle (1.8pt) node[above left] {$\gamma(t_4)$};
    \filldraw[black] (P5) circle (2pt) node[right] {$\gamma(t_5) = \gamma(b)$};

    \draw[crimsonred, <->] (P2) -- node[above, sloped] {\small Corda $\norm{\gamma(t_3)-\gamma(t_2)}$} (P3);

    \node[draw=crimsonred, fill=subtlegray, rounded corners, align=center] at (3.5,-0.3) 
        {\small Approssimazione poligonale: $L(\gamma) = \lim_{\abs{\mathcal{P}}\to 0} \sum \norm{\gamma(t_k) - \gamma(t_{k-1})} = \int_a^b \norm{\gamma'(t)}\dt$};
\end{tikzpicture}
\end{center}

\begin{teorema}{Rettificabilità e Formula Integrale della Lunghezza}{thm_lunghezza_integrale}
Ogni curva $\gamma: [a,b] \to \Rn$ di classe $C^1([a,b])$ (oppure regolare a tratti) è \textbf{rettificabile}, e la sua lunghezza è data dall'integrale di Riemann della norma del vettore velocità:
\[
L(\gamma) = \int_{a}^{b} \norm{\gamma'(t)} \dt = \int_{a}^{b} \sqrt{\sum_{i=1}^n [x_i'(t)]^2} \dt
\]
L'espressione differenziale scalare $\ds = \norm{\gamma'(t)}\dt$ è detta \textbf{elemento d'arco} (o metrica indotta).
\end{teorema}

\begin{dimostrazione}
Dimostriamo il teorema nel caso piano $n=2$ con $\gamma(t) = (x(t), y(t))$ di classe $C^1([a,b])$ (l'estensione a $n=3$ e a $\Rn$ è immediata).

Sia $\mathcal{P} = \graffe{a = t_0 < t_1 < \dots < t_N = b}$ una generica partizione di $[a,b]$. Consideriamo il generico incremento cordale su $[t_{k-1}, t_k]$:
\[
\Delta \gamma_k = \gamma(t_k) - \gamma(t_{k-1}) = \begin{pmatrix} x(t_k) - x(t_{k-1}) \\ y(t_k) - y(t_{k-1}) \end{pmatrix}
\]
Poiché $x(t)$ e $y(t)$ sono derivabili su ciascun intervallino $[t_{k-1}, t_k]$, per il \textbf{Teorema del Valor Medio di Lagrange} esistono due punti intermedi $\xi_k, \eta_k \in (t_{k-1}, t_k)$ tali che:
\[
x(t_k) - x(t_{k-1}) = x'(\xi_k)\Delta t_k, \quad y(t_k) - y(t_{k-1}) = y'(\eta_k)\Delta t_k
\]
dove $\Delta t_k = t_k - t_{k-1}$. La lunghezza della corda $k$-esima vale quindi:
\[
\norm{\Delta \gamma_k} = \sqrt{[x'(\xi_k)]^2 + [y'(\eta_k)]^2} \, \Delta t_k
\]
Sommando su tutti gli $N$ sottointervalli, la lunghezza poligonale è:
\[
L(\mathcal{P}, \gamma) = \sum_{k=1}^N \sqrt{[x'(\xi_k)]^2 + [y'(\eta_k)]^2} \, \Delta t_k
\]
Questa somma somiglia a una somma di Riemann per la funzione continua $g(t) = \norm{\gamma'(t)} = \sqrt{[x'(t)]^2 + [y'(t)]^2}$, con l'unica differenza che le derivate $x'$ e $y'$ sono valutate in punti generalmente distinti $\xi_k \neq \eta_k$.

Tuttavia, poiché $x'(t)$ e $y'(t)$ sono continue sull'intervallo compatto $[a,b]$, per il \textbf{Teorema di Heine-Cantor} esse sono \emph{uniformemente continue}: per ogni $\eps > 0$ esiste $\delta > 0$ tale che, se $\abs{\mathcal{P}} = \max_k \Delta t_k < \delta$, allora:
\[
\abs{\sqrt{[x'(\xi_k)]^2 + [y'(\eta_k)]^2} - \sqrt{[x'(\xi_k)]^2 + [y'(\xi_k)]^2}} < \eps
\]
Ne segue che la somma poligonale differisce da una vera somma di Riemann per una quantità infinitesima:
\[
\abs{L(\mathcal{P}, \gamma) - \sum_{k=1}^N \norm{\gamma'(\xi_k)}\Delta t_k} \le \sum_{k=1}^N \eps \, \Delta t_k = \eps (b - a)
\]
Passando al limite per il diametro della partizione $\abs{\mathcal{P}} \to 0$, la somma di Riemann converge all'integrale di Riemann della funzione continua $\norm{\gamma'(t)}$:
\[
\lim_{\abs{\mathcal{P}} \to 0} L(\mathcal{P}, \gamma) = \int_{a}^{b} \sqrt{[x'(t)]^2 + [y'(t)]^2} \dt = \int_{a}^{b} \norm{\gamma'(t)} \dt
\]
Poiché l'estremo superiore coincide con il limite dell'approssimazione poligonale, la tesi è pienamente dimostrata.
\end{dimostrazione}

\subsection{Formule Notevoli per la Lunghezza d'Arco}

A seconda del modo in cui una curva è assegnata, la formula generale $L = \int \norm{\gamma'(t)}\dt$ si specializza in espressioni operative di uso frequente:

\begin{metodo}[Formule Standard per la Lunghezza di Curve Notevoli]
\begin{enumerate}
    \item \textbf{Grafico Cartesiano di una Funzione $y = f(x)$ con $x \in [a,b]$}:
    Parametrizzazione naturale $\gamma(x) = (x, f(x)) \implies \gamma'(x) = (1, f'(x))$.
    \[
    \ds = \sqrt{1 + [f'(x)]^2} \dx \implies L = \int_{a}^{b} \sqrt{1 + [f'(x)]^2} \dx
    \]
    \item \textbf{Curva Piana in Coordinate Polari $r = \rho(\theta)$ con $\theta \in [\alpha, \beta]$}:
    Parametrizzazione cartesiana: $\gamma(\theta) = (\rho(\theta)\cos\theta, \rho(\theta)\sen\theta)$.\\
    Calcolando la derivata: $\gamma'(\theta) = (\rho'\cos\theta - \rho\sen\theta, \; \rho'\sen\theta + \rho\cos\theta)$.\\
    Quadrando le componenti e sfruttando l'identità $\cos^2\theta + \sen^2\theta = 1$:
    \[
    \norm{\gamma'(\theta)}^2 = (\rho'\cos\theta - \rho\sen\theta)^2 + (\rho'\sen\theta + \rho\cos\theta)^2 = [\rho(\theta)]^2 + [\rho'(\theta)]^2
    \]
    \[
    \ds = \sqrt{\rho(\theta)^2 + [\rho'(\theta)]^2} \dtheta \implies L = \int_{\alpha}^{\beta} \sqrt{\rho(\theta)^2 + [\rho'(\theta)]^2} \dtheta
    \]
    \item \textbf{Elica Cilindrica Circolare in $\Rtre$}:
    Parametrizzazione $\gamma(t) = (R\cos t, R\sen t, c t)$ con $t \in [0, 2\pi N]$ (avente passo dell'elica $h = 2\pi c$ e raggio $R$).
    \begin{align*}
    \gamma'(t) &= (-R\sen t, R\cos t, c) \\
    \implies \norm{\gamma'(t)} &= \sqrt{R^2\sen^2 t + R^2\cos^2 t + c^2} = \sqrt{R^2 + c^2}
    \end{align*}
    La velocità scalare è \textbf{costante}, per cui la lunghezza per un singolo giro ($t \in [0, 2\pi]$) è semplicemente:
    \[
    L = \int_{0}^{2\pi} \sqrt{R^2 + c^2} \dt = 2\pi \sqrt{R^2 + c^2}
    \]
\end{enumerate}
\end{metodo}

\begin{teorema}{La Retta come Curva di Minima Lunghezza (Geodesica Euclidea)}{geodesica_euclidea_completa}
Siano $P, Q \in \Rn$ due punti distinti. Tra tutte le curve $C^1$ regolari a tratti $\gamma: [a,b] \to \Rn$ che congiungono $P$ e $Q$ (ossia tali che $\gamma(a) = P$ e $\gamma(b) = Q$), il cammino di lunghezza minima è il \textbf{segmento rettilineo} orientato da $P$ a $Q$.
\end{teorema}

\begin{dimostrazione}
Sia $\mathbf{u} = \frac{Q - P}{\norm{Q - P}}$ il versore unitario che individua la direzione e il verso da $P$ a $Q$ (dunque $\norm{\mathbf{u}} = 1$).\\
Poiché $\gamma$ è di classe $C^1$ a tratti con $\gamma(a) = P$ e $\gamma(b) = Q$, per il Teorema Fondamentale del Calcolo Integrale possiamo esprimere il vettore spostamento totale come l'integrale del vettore velocità:
\[
Q - P = \gamma(b) - \gamma(a) = \int_{a}^{b} \gamma'(t) \dt
\]
Calcoliamo ora il prodotto scalare tra il vettore spostamento $Q - P$ e il versore $\mathbf{u}$:
\[
\scal{Q - P, \mathbf{u}} = \scal{\int_{a}^{b} \gamma'(t) \dt, \mathbf{u}} = \int_{a}^{b} \scal{\gamma'(t), \mathbf{u}} \dt
\]
D'altra parte, per definizione di $\mathbf{u}$, il membro di sinistra è esattamente la distanza euclidea tra i due punti:
\[
\scal{Q - P, \mathbf{u}} = \scal{Q - P, \frac{Q - P}{\norm{Q - P}}} = \frac{\norm{Q - P}^2}{\norm{Q - P}} = \norm{Q - P}
\]
Applichiamo la \textbf{disuguaglianza di Cauchy-Schwarz} al prodotto scalare integrando:
\[
\scal{\gamma'(t), \mathbf{u}} \le \abs{\scal{\gamma'(t), \mathbf{u}}} \le \norm{\gamma'(t)} \norm{\mathbf{u}} = \norm{\gamma'(t)} \cdot 1 = \norm{\gamma'(t)}
\]
Integrando ambo i membri sull'intervallo $[a,b]$:
\[
\norm{Q - P} = \int_{a}^{b} \scal{\gamma'(t), \mathbf{u}} \dt \le \int_{a}^{b} \norm{\gamma'(t)} \dt = L(\gamma)
\]
Abbiamo così dimostrato che la lunghezza $L(\gamma)$ di una qualsiasi curva congiungente $P$ e $Q$ è sempre maggiore o uguale alla distanza euclidea $\norm{Q - P}$:
\[
L(\gamma) \ge \norm{Q - P}
\]
L'uguaglianza $L(\gamma) = \norm{Q - P}$ sussiste se e solo se la disuguaglianza di Cauchy-Schwarz è soddisfatta come uguaglianza quasi ovunque su $[a,b]$ con segno positivo:
\[
\scal{\gamma'(t), \mathbf{u}} = \norm{\gamma'(t)} \iff \gamma'(t) = \lambda(t) \mathbf{u}, \quad \text{con } \lambda(t) = \norm{\gamma'(t)} > 0
\]
Ciò implica che il vettore velocità $\gamma'(t)$ deve mantenere costantemente la stessa direzione e lo stesso verso del versore $\mathbf{u}$. Integrando rispetto al tempo:
\[
\gamma(t) = \gamma(a) + \left(\int_a^t \lambda(u)\du\right) \mathbf{u} = P + s(t) \frac{Q - P}{\norm{Q - P}}
\]
che è precisamente la parametrizzazione del \textbf{segmento di retta} che unisce $P$ a $Q$.
\end{dimostrazione}

% ==============================================================================
\section{Integrali Curvilinei di Prima Specie (di Campi Scalari)}
% ==============================================================================

Gli integrali curvilinei di prima specie consentono di integrare una funzione scalare $f: \Rn \to \R$ lungo una curva unidimensionale. 

\subsection{Motivazione Fisica e Geometrica}
\begin{enumerate}
    \item \textbf{Massa di un filo materiale}: Se immaginiamo che la curva $\Gamma \subset \Rtre$ rappresenti un filo sottile la cui densità lineare di massa varia da punto a punto secondo una legge continua $\rho(x,y,z)$ $[\text{kg/m}]$, la massa totale del filo si ottiene sommando i contributi elementari $\diff M = \rho(\gamma(t))\ds$:
    \[
    M = \int_\gamma \rho(x,y,z) \ds
    \]
    \item \textbf{Baricentro e Momenti d'Inerzia}: Le coordinate del centro di massa $G = (x_G, y_G, z_G)$ del filo sono le medie ponderate delle coordinate rispetto alla massa:
    \[
    x_G = \frac{1}{M} \int_\gamma x \, \rho \ds, \quad y_G = \frac{1}{M} \int_\gamma y \, \rho \ds, \quad z_G = \frac{1}{M} \int_\gamma z \, \rho \ds
    \]
    mentre il momento d'inerzia rispetto all'asse $z$ vale $I_z = \int_\gamma (x^2 + y^2) \rho \ds$.
    \item \textbf{Area di una Cortina/Staccionata Cilindrica}: Se $\gamma$ è una curva piana nel piano $xy$ e $f(x,y) \ge 0$ è una funzione quota, l'integrale $\int_\gamma f \ds$ rappresenta l'area della superficie verticale (la "staccionata") che ha per base il sostegno $\Gamma$ sul piano $z=0$ e altezza variabile pari a $z = f(x,y)$.
\end{enumerate}

\begin{definizione}{Integrale Curvilineo di I Specie}{def_integrale_I_specie}
Sia $\gamma: [a,b] \to \Rn$ una curva regolare (o regolare a tratti) con sostegno $\Gamma = \gamma([a,b])$, e sia $f: \Gamma \to \R$ una funzione scalare continua. Si definisce \textbf{integrale curvilineo di prima specie} di $f$ lungo $\gamma$:
\[
\int_\gamma f \ds = \int_{\gamma} f(\bx) \ds = \int_{a}^{b} f(\gamma(t)) \, \norm{\gamma'(t)} \dt
\]
Nel piano $\Rdue$, se $\gamma(t) = (x(t), y(t))$:
\[
\int_\gamma f(x,y) \ds = \int_{a}^{b} f(x(t), y(t)) \sqrt{[x'(t)]^2 + [y'(t)]^2} \dt
\]
\end{definizione}

\begin{center}
\begin{tikzpicture}[scale=1.1, >=Stealth]
    % Assi 3D per l'interpretazione geometrica della staccionata
    \draw[->, thick, gray] (0,0) -- (6,0) node[right, black] {$y$};
    \draw[->, thick, gray] (0,0) -- (-2,-1.5) node[below left, black] {$x$};
    \draw[->, thick, gray] (0,0) -- (0,4) node[above, black] {$z$};

    % Curva di base nel piano xy
    \draw[ultra thick, navyblue] (1,-0.5) to[out=20, in=-120] (3.5,-0.2) to[out=60, in=-150] (5,1.2);
    \node[navyblue, below right] at (5,1.2) {Base $\Gamma = \gamma(t)$};

    % Curva superiore z = f(x,y)
    \draw[ultra thick, crimsonred] (1,1.8) to[out=20, in=-120] (3.5,2.6) to[out=60, in=-150] (5,3.2);
    \node[crimsonred, above] at (3.5,2.7) {$z = f(x,y)$};

    % Righe verticali per formare la staccionata
    \foreach \x/\y/\z in {1/-0.5/1.8, 1.6/-0.42/2.1, 2.2/-0.35/2.35, 2.8/-0.28/2.52, 3.5/-0.2/2.6, 4.1/0.3/2.8, 4.6/0.8/3.05, 5/1.2/3.2} {
        \draw[thin, coolblue!60] (\x,\y) -- (\x,\z);
    }

    % Riempimento sfumato staccionata
    \shade[top color=coolblue!30, bottom color=coolblue!5, opacity=0.7] 
        (1,-0.5) to[out=20, in=-120] (3.5,-0.2) to[out=60, in=-150] (5,1.2) --
        (5,3.2) to[out=-150, in=60] (3.5,2.6) to[out=-120, in=20] (1,1.8) -- cycle;

    \node[draw=coolblue!80!black, fill=white, rounded corners, align=center] at (2.2,3.4) 
        {\small \textbf{Area Staccionata}: $\text{Area} = \int_\gamma f(x,y) \ds$};
\end{tikzpicture}
\end{center}

\begin{teorema}{Invarianza per Riparametrizzazione e per Inversione di Verso}{thm_invarianza_I_specie}
L'integrale curvilineo di prima specie gode delle seguenti proprietà fondamentali:
\begin{enumerate}
    \item \textbf{Invarianza rispetto a qualsiasi riparametrizzazione ammissibile}: Se $\tilde{\gamma}(s) = \gamma(\phi(s))$ con $\phi: [c,d] \to [a,b]$ diffeomorfismo $C^1$, allora:
    \[
    \int_{\tilde{\gamma}} f \ds = \int_\gamma f \ds
    \]
    \item \textbf{Invarianza rispetto al verso di percorrenza}: Se indichiamo con $\gamma^-$ la curva percorsa nel verso opposto, definita da $\gamma^-(t) = \gamma(a + b - t)$ per $t \in [a,b]$, si ha:
    \[
    \int_{\gamma^-} f \ds = \int_\gamma f \ds
    \]
\end{enumerate}
\end{teorema}

\begin{dimostrazione}
Dimostriamo innanzitutto la proprietà generale (1). Sia $t = \phi(s)$ con $\phi: [c,d] \to [a,b]$.
Calcoliamo l'integrale lungo la curva riparametrizzata $\tilde{\gamma}$:
\[
\int_{\tilde{\gamma}} f \ds = \int_{c}^{d} f(\tilde{\gamma}(s)) \norm{\tilde{\gamma}'(s)} \ds = \int_{c}^{d} f(\gamma(\phi(s))) \norm{\gamma'(\phi(s)) \, \phi'(s)} \ds
\]
Sfruttando l'omogeneità della norma rispetto agli scalari reali, $\norm{\gamma'(\phi(s)) \, \phi'(s)} = \norm{\gamma'(\phi(s))} \abs{\phi'(s)}$. Dunque:
\[
\int_{\tilde{\gamma}} f \ds = \int_{c}^{d} f(\gamma(\phi(s))) \norm{\gamma'(\phi(s))} \abs{\phi'(s)} \ds
\]
Distinguiamo i due casi possibili:
\begin{itemize}
    \item Se $\phi'(s) > 0$ (riparametrizzazione concorde), allora $\abs{\phi'(s)} = \phi'(s)$, $\phi(c)=a$, $\phi(d)=b$. Effettuando la sostituzione di variabile $t = \phi(s)$ con $\dt = \phi'(s)\ds$:
    \[
    \int_{c}^{d} f(\gamma(\phi(s))) \norm{\gamma'(\phi(s))} \phi'(s) \ds = \int_{a}^{b} f(\gamma(t)) \norm{\gamma'(t)} \dt = \int_\gamma f \ds
    \]
    \item Se $\phi'(s) < 0$ (riparametrizzazione discorde), allora $\abs{\phi'(s)} = -\phi'(s)$, $\phi(c)=b$, $\phi(d)=a$. Sostituendo $t = \phi(s)$ e invertendo gli estremi di integrazione:
    \begin{align*}
    \int_{c}^{d} f(\gamma(\phi(s))) \norm{\gamma'(\phi(s))} (-\phi'(s)) \ds &= -\int_{b}^{a} f(\gamma(t)) \norm{\gamma'(t)} \dt \\
    &= \int_{a}^{b} f(\gamma(t)) \norm{\gamma'(t)} \dt = \int_\gamma f \ds
    \end{align*}
\end{itemize}
La proprietà (2) è un caso particolare di riparametrizzazione discorde con $\phi(t) = a + b - t$, per cui $\phi'(t) = -1$ e $\abs{\phi'(t)} = 1$. L'integrale non cambia segno.
\end{dimostrazione}

\begin{esempio}{Calcolo del Baricentro di un Filo Semicircolare}{es_baricentro_semicircolare}
Calcolare la posizione del baricentro $G = (x_G, y_G)$ di un filo semicircolare omogeneo di raggio $R$ posto nel semipiano superiore:
\[
\Gamma = \graffe{(x,y) \in \Rdue : x^2 + y^2 = R^2, \; y \ge 0}
\]
avente densità lineare di massa costante $\rho_0 > 0$.

\textbf{Risoluzione passo-passo}:
\begin{enumerate}
    \item \textbf{Parametrizzazione della curva}:
    \[
    \gamma(t) = \begin{pmatrix} R\cos t \\ R\sen t \end{pmatrix}, \quad t \in [0, \pi]
    \]
    Il vettore velocità e l'elemento d'arco scalare sono:
    \[
    \gamma'(t) = \begin{pmatrix} -R\sen t \\ R\cos t \end{pmatrix} \implies \norm{\gamma'(t)} = \sqrt{R^2\sen^2 t + R^2\cos^2 t} = R \implies \ds = R\dt
    \]
    \item \textbf{Massa totale del filo}:
    \[
    M = \int_\gamma \rho_0 \ds = \rho_0 \int_{0}^{\pi} R \dt = \pi R \rho_0
    \]
    \item \textbf{Ascissa del baricentro $x_G$}:
    Per ragioni di simmetria rispetto all'asse $y$ (densità costante e dominio pari in $x$), ci aspettiamo $x_G = 0$. Verifichiamo con il calcolo analitico:
    \[
    x_G = \frac{1}{M} \int_\gamma x \, \rho_0 \ds = \frac{\rho_0}{\pi R \rho_0} \int_{0}^{\pi} (R\cos t) R \dt = \frac{R}{\pi} \left[ \sen t \right]_{0}^{\pi} = \frac{R}{\pi} (0 - 0) = 0
    \]
    \item \textbf{Ordinata del baricentro $y_G$}:
    \begin{align*}
    y_G &= \frac{1}{M} \int_\gamma y \, \rho_0 \ds = \frac{\rho_0}{\pi R \rho_0} \int_{0}^{\pi} (R\sen t) R \dt \\
    &= \frac{R}{\pi} \int_{0}^{\pi} \sen t \dt = \frac{R}{\pi} \left[ -\cos t \right]_{0}^{\pi} = \frac{R}{\pi} (1 - (-1)) = \frac{2R}{\pi}
    \end{align*}
\end{enumerate}
Il baricentro si trova dunque nel punto $G = \left(0, \frac{2R}{\pi}\right) \approx (0, 0.6366 R)$, che cade rigorosamente all'esterno del sostegno materiale della curva.
\end{esempio}

% ==============================================================================
\section{Integrali Curvilinei di Seconda Specie (Lavoro e Circolazione)}
% ==============================================================================

Mentre gli integrali di prima specie associano a ogni elemento di curva uno scalare positivo $\ds = \norm{\gamma'(t)}\dt$, gli \textbf{integrali di seconda specie} integrano un \textbf{campo vettoriale} lungo il \textbf{vettore spostamento orientato} $\diff\mathbf{r} = \gamma'(t)\dt$.

\subsection{Motivazione Fisica: Il Lavoro di una Forza}

In fisica classica, il lavoro elementare $\diff W$ compiuto da una forza costante $\bF$ per uno spostamento rettilineo $\Delta \br$ è dato dal prodotto scalare $\diff W = \bF \cdot \Delta \br = \norm{\bF} \norm{\Delta \br} \cos\theta$.

Quando una particella si muove lungo una traiettoria curvilinea $\gamma(t)$ soggetta a un campo di forze variabile nello spazio $\bF(\bx) = (F_1(\bx), \dots, F_n(\bx))$, la potenza istantanea erogata dalla forza al tempo $t$ è il prodotto scalare tra la forza e la velocità:
\[
P(t) = \scal{\bF(\gamma(t)), \gamma'(t)}
\]
Il lavoro complessivo compiuto dal campo lungo l'intero percorso temporale $[a,b]$ è l'integrale della potenza nel tempo:
\[
W = \int_{a}^{b} P(t) \dt = \int_{a}^{b} \scal{\bF(\gamma(t)), \gamma'(t)} \dt
\]

\begin{definizione}{Integrale Curvilineo di II Specie}{def_integrale_II_specie}
Sia $A \subseteq \Rn$ un aperto, $\bF: A \to \Rn$ un campo vettoriale continuo ($\bF = (F_1, \dots, F_n)$) e sia $\gamma: [a,b] \to A$ una curva orientata regolare a tratti. Si definisce \textbf{integrale curvilineo di seconda specie} (o \textbf{lavoro}) di $\bF$ lungo $\gamma$:
\[
\int_\gamma \bF \cdot \diff\mathbf{r} = \int_\gamma \sum_{i=1}^n F_i \dx_i = \int_{a}^{b} \scal{\mathbf{F}(\gamma(t)), \gamma'(t)} \dt = \int_{a}^{b} \left( \sum_{i=1}^n F_i(\gamma(t)) \, x_i'(t) \right) \dt
\]
Se la curva $\gamma$ è chiusa ($\gamma(a) = \gamma(b)$), l'integrale è detto \textbf{circolazione} (o \emph{circuitazione}) del campo lungo $\gamma$ e viene indicato con il simbolo:
\[
\oint_\gamma \bF \cdot \diff\mathbf{r}
\]
\end{definizione}

\begin{center}
\begin{tikzpicture}[scale=1.2, >=Stealth]
    % Traiettoria orientata
    \draw[ultra thick, navyblue, ->] (0,0) to[out=30, in=-140] (3,1.5) to[out=40, in=-160] (6,2.2);
    \node[navyblue, below right] at (6,2.2) {$\gamma(t)$};

    % Punto generico sulla curva
    \coordinate (P) at (3,1.5);
    \filldraw[black] (P) circle (2pt) node[below=3pt] {$\gamma(t)$};

    % Vettore spostamento dr = gamma'(t) dt
    \draw[thick, crimsonred, ->, line width=1.3pt] (P) -- ++(1.4, 0.7) node[below right] {$\diff\mathbf{r} = \gamma'(t)\dt$};

    % Campo di forze F(gamma(t))
    \draw[thick, forestgreen, ->, line width=1.4pt] (P) -- ++(0.8, 1.6) node[above right] {$\mathbf{F}(\gamma(t))$};

    % Angolo theta tra F e dr
    \draw[thin, black] (3.6,1.8) arc (30:63:0.7);
    \node at (3.8,2.05) {\small $\theta$};

    % Proiezione tratteggiata di F lungo dr
    \draw[dashed, gray] (3.8, 3.1) -- (4.4, 2.2);
    \draw[thick, warmamber!90!black, ->] (P) -- (4.4, 2.2) node[midway, below] {\small $F_T = \norm{\bF}\cos\theta$};

    \node[draw=navyblue, fill=subtlegray, rounded corners, align=center] at (3,-0.6) 
        {\small Lavoro elementare: $\diff W = \scal{\mathbf{F}, \diff\mathbf{r}} = \norm{\mathbf{F}} \norm{\diff\mathbf{r}} \cos\theta = \scal{\mathbf{F}(\gamma(t)), \gamma'(t)}\dt$};
\end{tikzpicture}
\end{center}

\subsection{Comportamento Rispetto al Cambio di Verso (Antimetria)}

La differenza più profonda tra gli integrali di I e di II specie risiede nel comportamento rispetto all'inversione del verso di percorrenza.

\begin{teorema}{Antimetria degli Integrali di II Specie}{thm_antimetria_II_specie}
Sia $\gamma: [a,b] \to A$ una curva regolare a tratti, e sia $\gamma^-(t) = \gamma(a + b - t)$ la curva percorsa in verso opposto. Allora:
\[
\int_{\gamma^-} \bF \cdot \diff\mathbf{r} = -\int_{\gamma} \bF \cdot \diff\mathbf{r}
\]
\end{teorema}

\begin{dimostrazione}
Sia $\gamma^-(t) = \gamma(\phi(t))$ con $\phi(t) = a + b - t$ per $t \in [a,b]$.
La derivata del cambio di variabile è $\phi'(t) = -1$, per cui la velocità della curva invertita è:
\[
(\gamma^-)'(t) = \der{}{t}\left[\gamma(a+b-t)\right] = -\gamma'(a+b-t)
\]
Applicando la definizione di integrale di seconda specie lungo $\gamma^-$:
\[
\int_{\gamma^-} \bF \cdot \diff\mathbf{r} = \int_{a}^{b} \scal{\bF(\gamma^-(t)), (\gamma^-)'(t)} \dt = \int_{a}^{b} \scal{\bF(\gamma(a+b-t)), -\gamma'(a+b-t)} \dt
\]
Portando il segno meno fuori dal prodotto scalare ed effettuando il cambio di variabile $\tau = a + b - t$ con $\diff\tau = -\dt$:
\begin{align*}
\int_{\gamma^-} \bF \cdot \diff\mathbf{r} &= -\int_{a}^{b} \scal{\bF(\gamma(a+b-t)), \gamma'(a+b-t)} \dt \\
&= -\int_{b}^{a} \scal{\bF(\gamma(\tau)), \gamma'(\tau)} (-\diff\tau) \\
&= \int_{b}^{a} \scal{\bF(\gamma(\tau)), \gamma'(\tau)} \diff\tau \\
&= -\int_{a}^{b} \scal{\bF(\gamma(\tau)), \gamma'(\tau)} \diff\tau = -\int_{\gamma} \bF \cdot \diff\mathbf{r}
\end{align*}
Il cambio di segno è dovuto al fatto che il vettore spostamento istantaneo inverte la propria direzione spaziale ($\diff\mathbf{r}^- = -\diff\mathbf{r}$).
\end{dimostrazione}

\begin{esame}[Confronto Chiave tra Integrali di I Specie e II Specie]
Per evitare gli errori più gravi nei compiti scritti e nei quesiti orali, si tenga sempre a mente la seguente tavola sinottica:
\begin{center}
\begin{tabularx}{\textwidth}{lXX}
\toprule
\textbf{Caratteristica} & \textbf{Integrale di I Specie $\int_\gamma f \ds$} & \textbf{Integrale di II Specie $\int_\gamma \bF \cdot \diff\br$} \\
\midrule
\textbf{Oggetto integrando} & Campo scalare $f: \Rn \to \R$ & Campo vettoriale $\bF: \Rn \to \Rn$ \\
\textbf{Elemento differenziale} & Scalare positivo: $\ds = \norm{\gamma'(t)}\dt$ & Vettore orientato: $\diff\br = \gamma'(t)\dt$ \\
\textbf{Formula di calcolo} & $\int_a^b f(\gamma(t)) \norm{\gamma'(t)}\dt$ & $\int_a^b \scal{\bF(\gamma(t)), \gamma'(t)}\dt$ \\
\textbf{Inversione di verso} & \textbf{Invariante}: $\int_{\gamma^-} f\ds = +\int_\gamma f\ds$ & \textbf{Antimetrico}: $\int_{\gamma^-} \bF\cdot\diff\br = -\int_\gamma \bF\cdot\diff\br$ \\
\textbf{Significato fisico tipico} & Massa di un filo, area di una cortina & Lavoro di una forza, circuitazione \\
\bottomrule
\end{tabularx}
\end{center}
\end{esame}

\begin{esempio}{Calcolo del Lavoro su Cammini Differenti}{es_lavoro_cammini_diversi}
Si consideri il campo vettoriale piano $\bF(x,y) = (2xy, \; x^2 - y^2)$ definito su tutto $\Rdue$. Calcolare il lavoro compiuto da $\bF$ per spostare una particella dall'origine $O = (0,0)$ al punto $P = (1,1)$ lungo due percorsi distinti:
\begin{enumerate}
    \item Il segmento rettilineo $\gamma_1$ congiungente $(0,0)$ a $(1,1)$;
    \item L'arco di parabola $\gamma_2$ di equazione $y = x^2$ con $x \in [0,1]$.
\end{enumerate}

\textbf{Svolgimento}:
\begin{enumerate}
    \item \textbf{Lungo il segmento $\gamma_1$}:
    Parametrizziamo il segmento ponendo $\gamma_1(t) = (t, t)$ con $t \in [0, 1]$.
    Il vettore velocità è $\gamma_1'(t) = (1, 1)$.
    Sostituendo le componenti lungo la curva:
    \[
    \bF(\gamma_1(t)) = \begin{pmatrix} 2(t)(t) \\ t^2 - t^2 \end{pmatrix} = \begin{pmatrix} 2t^2 \\ 0 \end{pmatrix}
    \]
    Il lavoro lungo $\gamma_1$ vale:
    \[
    W_1 = \int_{\gamma_1} \bF \cdot \diff\mathbf{r} = \int_{0}^{1} \scal{\begin{pmatrix} 2t^2 \\ 0 \end{pmatrix}, \begin{pmatrix} 1 \\ 1 \end{pmatrix}} \dt = \int_{0}^{1} 2t^2 \dt = \left[ \frac{2}{3} t^3 \right]_{0}^{1} = \frac{2}{3}
    \]

    \item \textbf{Lungo l'arco di parabola $\gamma_2$}:
    Parametrizziamo la parabola ponendo $\gamma_2(t) = (t, t^2)$ con $t \in [0, 1]$.
    Il vettore velocità è $\gamma_2'(t) = (1, 2t)$.
    Valutando il campo sulla parabola:
    \[
    \bF(\gamma_2(t)) = \begin{pmatrix} 2(t)(t^2) \\ t^2 - (t^2)^2 \end{pmatrix} = \begin{pmatrix} 2t^3 \\ t^2 - t^4 \end{pmatrix}
    \]
    Il prodotto scalare tra il campo e il vettore velocità è:
    \[
    \scal{\bF(\gamma_2(t)), \gamma_2'(t)} = (2t^3)(1) + (t^2 - t^4)(2t) = 2t^3 + 2t^3 - 2t^5 = 4t^3 - 2t^5
    \]
    Integrando rispetto a $t \in [0, 1]$:
    \[
    W_2 = \int_{\gamma_2} \bF \cdot \diff\mathbf{r} = \int_{0}^{1} (4t^3 - 2t^5) \dt = \left[ t^4 - \frac{t^6}{3} \right]_{0}^{1} = 1 - \frac{1}{3} = \frac{2}{3}
    \]
\end{enumerate}
I due integrali coincidono esattamente ($W_1 = W_2 = 2/3$). È una pura coincidenza numerica o nasconde una proprietà strutturale profonda del campo vettoriale $\bF$? Come scopriremo nel prossimo capitolo, $\bF$ è un \textbf{campo conservativo}, il cui lavoro non dipende dalla forma della traiettoria ma unicamente dai punti di partenza e di arrivo.
\end{esempio}
